\subsection{Characterization of the nilpotent elements of a ring}

PROPOSITION 1. --- Let A be a commutative ring and x an element of A. For x to be nilpotent it is necessary and sufficient that 1 - xT be invertible in the ring A{[}T{]}.

We note that A{[}T{]} is a subring of the formal power series ring A{[}{[}T{]}{]} and that 1 - xT has in A{[}{[}T{]}{]} the inverse \(\sum_{n=0}^{\infty} x^{n} T^{n}\) (IV, p.~30, Prop. 5). For 1 - xT to be

invertible in A{[}T{]} it is necessary and sufficient for \(\sum_{n=0}^{\infty} x^n T^n\) to be a polynomial,

that is, for \(x\) to be nilpotent.

PROPOSITION 2. --- Let A be a commutative ring. The set of nilpotent elements of A is an ideal of A equal to the intersection of the set of all prime ideals of A.

Let \(x\) be a nilpotent element of \(\mathbf{A}\) and \(\mathfrak{p}\) a prime ideal of \(\mathbf{A}\) . The residue class of \(x \mod \mathfrak{p}\) is a nilpotent element of the integral domain \(\mathbf{A} / \mathfrak{p}\) , hence is zero; thus we have \(x \in \mathfrak{p}\) .

Let \(x\) be a non-nilpotent element of \(A\) . By Prop. 1 the principal ideal \((1 - xT)\) of \(A[T]\) is distinct from \(A[T]\) . By Krull's theorem (I, p.~104) there exists a maximal ideal \(m\) of \(A[T]\) containing \(1 - xT\) . Then \(m\) is a prime ideal of \(A[T]\) , hence \(p = A \cap m\) is a prime ideal of \(A\) . We have \(1 \notin m\) and \(1 - Tx \in m\) , hence \(Tx \notin m\) and a fortiori \(x \notin p\) .

We have thus shown that the set n of nilpotent elements of A is the intersection of the set of all prime ideals of A; since every intersection of ideals is an ideal, n is an ideal.

COROLLARY. --- For a commutative ring to be reduced (V, p.~34, Def. 2) it is necessary and sufficient for it to be isomorphic to a subring of a product of fields.

The condition is clearly sufficient.

Let A be reduced. By Prop. 2 the intersection n of the set of prime ideals of A is reduced to 0. For every prime ideal p of A, let \(k(\mathfrak{p})\) be the field of fractions of A/p and \(\varphi_{p}\) the canonical homomorphism of A into \(k(\mathfrak{p})\) . Let \(\varphi\) be the homomorphism of A into \(\prod_{\mathfrak{p}} k(\mathfrak{p})\) whose component of index p is \(\varphi_{p}\) . The kernel of \(\varphi_{\mathfrak{p}}\) is \(\mathfrak{p}\) , hence that of \(\varphi\) is \(\mathfrak{n} = 0\) ; therefore \(\varphi\) is an isomorphism of \(\mathbf{A}\) onto a subring of \(\prod_{\mathfrak{p}} k(\mathfrak{p})\) .

